% Standalone problem document, generated from the adjacent README.md.
% Compile with XeLaTeX. Mathematical content is maintained in Markdown.
\documentclass[11pt,a4paper]{article}
\usepackage[fontsize=10.5pt]{scrextend}
\usepackage[margin=22mm,headheight=15pt,headsep=9mm,footskip=11mm]{geometry}
\usepackage{fontspec}
\setmainfont{texgyrepagella}[Extension=.otf,UprightFont=*-regular,BoldFont=*-bold,ItalicFont=*-italic,BoldItalicFont=*-bolditalic]
\setsansfont{texgyreheros}[Extension=.otf,UprightFont=*-regular,BoldFont=*-bold,ItalicFont=*-italic,BoldItalicFont=*-bolditalic]
\setmonofont{texgyrecursor}[Extension=.otf,UprightFont=*-regular,BoldFont=*-bold,ItalicFont=*-italic,BoldItalicFont=*-bolditalic,Scale=MatchLowercase]
\usepackage{amsmath,amssymb}
\usepackage{unicode-math}
\setmathfont{texgyrepagella-math.otf}
\usepackage{xcolor,microtype,enumitem,needspace,fancyhdr}
\usepackage{bookmark}
\definecolor{ink}{HTML}{17344B}
\definecolor{accent}{HTML}{197D80}
\definecolor{muted}{HTML}{596773}
\hypersetup{colorlinks=true,linkcolor=accent,urlcolor=accent,pdftitle={MI-04 - A
universal block-norm characterization of essentially Hermitian
matrices},pdfauthor={Open Problems in Numerical Linear Algebra}}
\urlstyle{same}
\setlength{\parindent}{0pt}
\setlength{\parskip}{5pt plus 1pt minus 1pt}
\linespread{1.04}
\setlength{\emergencystretch}{3em}
\widowpenalty=10000
\clubpenalty=10000
\displaywidowpenalty=10000
\allowdisplaybreaks[1]
\setlist{nosep,leftmargin=1.5em}
\providecommand{\tightlist}{\setlength{\itemsep}{0pt}\setlength{\parskip}{0pt}}
\providecommand{\pandocbounded}[1]{#1}
\pagestyle{fancy}
\fancyhf{}
\fancyhead[L]{\sffamily\footnotesize\color{muted}OPEN PROBLEMS IN NUMERICAL LINEAR ALGEBRA}
\fancyhead[R]{\sffamily\bfseries\footnotesize\color{accent}MI-04}
\fancyfoot[L]{\parbox[b]{0.9\textwidth}{\sffamily\fontsize{8}{10}\selectfont\color{muted}Editorial ratings; a bounded search cannot certify that a problem remains unsolved.\\Literature check: 2026-09-10}}
\fancyfoot[R]{\sffamily\footnotesize\color{muted}\thepage}
\renewcommand{\headrulewidth}{0.3pt}
\renewcommand{\headrule}{\hbox to\headwidth{\color{accent}\leaders\hrule height \headrulewidth\hfill}}
\makeatletter
\renewcommand{\section}{\Needspace{5\baselineskip}\@startsection{section}{1}{0pt}{12pt plus 2pt minus 2pt}{4pt}{\sffamily\bfseries\large\color{ink}}}
\renewcommand{\subsection}{\Needspace{4\baselineskip}\@startsection{subsection}{2}{0pt}{9pt plus 1pt minus 1pt}{3pt}{\sffamily\bfseries\normalsize\color{ink}}}
\makeatother
\setcounter{secnumdepth}{0}
\begin{document}
{\sffamily\small\color{muted}Matrix inequalities and norms\par}
\vspace{3pt}
{\raggedright\hyphenpenalty=10000\exhyphenpenalty=10000\sffamily\bfseries\fontsize{21}{24}\selectfont\color{ink}A
universal block-norm characterization of essentially Hermitian
matrices\par}
\vspace{7pt}
{\sffamily\small\textbf{Difficulty:} challenging\quad\textbf{Importance:} interesting
to specialist\par}
{\sffamily\small\bfseries\color{accent}Status: Open\par}
\vspace{4pt}
{\color{accent}\hrule height 0.8pt}
\vspace{6pt}
\textbf{Rating rationale:} The universal converse needs new control of
positive block completions; its immediate impact is a specific
numerical-range characterization.

\subsection{Problem statement}\label{problem-statement}

Let \(n\ge1\) and \(X\in\mathbb C^{n\times n}\). Suppose that for every
pair of Hermitian \(A,B\in\mathbb C^{n\times n}\) for which

\[H=\begin{bmatrix}A&X\\X^*&B\end{bmatrix}\succeq0,\]

one has \(\|H\|_2\le\|A+B\|_2\), where \(\|\cdot\|_2\) is the operator
norm. Must there exist a Hermitian \(K\) and scalars
\(\alpha,\beta\in\mathbb C\) with \(X=\alpha K+\beta I_n\)?

Such an \(X\) is called essentially Hermitian; equivalently its
numerical range \(\{v^*Xv:\|v\|_2=1\}\) lies in an affine line.

\subsection{Why it matters}\label{why-it-matters}

The question characterizes exactly when off-diagonal coupling in a
positive block matrix can always be ignored in a particular
spectral-norm bound. It connects a numerical-range geometry condition to
a universal norm estimate.

\subsection{References}\label{references}

\begin{enumerate}
\def\labelenumi{\arabic{enumi}.}
\tightlist
\item
  J.-C. Bourin and E.-Y. Lee, \emph{Eigenvalue inequalities for positive
  block matrices with the inradius of the numerical range},
  arXiv:2111.15180v1 (30 November 2021), Conjecture 3.3, Theorem 3.2,
  and Proposition 3.4. \href{https://arxiv.org/html/2111.15180}{Primary
  text}.
\item
  T. Hayashi, \emph{On a norm inequality for a positive block-matrix},
  Linear Algebra and its Applications 566 (2019), 86--97, Theorem 2.5.
  \href{https://arxiv.org/abs/1808.00181}{Preprint},
  \href{https://doi.org/10.1016/j.laa.2018.12.027}{DOI}.
\end{enumerate}

\subsection{Status check --- 2026-09-10}\label{status-check-2026-09-10}

The latest arXiv version of reference 1 remains v1. Hayashi proves
normality under additional invertibility and distinct-singular-value
assumptions, not the stated conclusion in full generality. A
Frobenius-norm variant characterizes normality and is already proved.
Searches included \textit{Bourin essentially Hermitian conjecture},
\textit{universal positive block matrix norm Hayashi conjecture},
and \textit{essentially Hermitian conjecture proof 2025 2026}. No
later resolution of Conjecture 3.3 was located. A 2025 result on
decomposable numerical ranges settles a different conjecture.

\textbf{Audit update (2026-09-10):} Rechecked Bourin--Lee Conjecture 3.3
and searched for subsequent essentially-Hermitian characterizations. The
Frobenius-norm theorem has a different norm and conclusion; no
resolution of the operator-norm converse was located. This is a bounded
literature check, not a proof that no solution exists.
\end{document}
